\section{Relatório Final -- SIGAA}
\label{cap:Relatorio_Final_SIGAA}

Relatório submetido no SIGAA ao final do projeto.
Apenas texto simples pode ser utilizado para responder a estes campos.

\subsection{Resumo}

\subsubsection{Resumo}

Este trabalho investiga a implementação de métodos numéricos em diferentes
paradigmas de programação, com foco na comparação entre uma linguagem
multiparadigma, Python, e uma linguagem funcional pura, Haskell. Foram
implementados métodos clássicos de cálculo numérico, incluindo o método do
ponto fixo, o método de Newton--Raphson e o método de Euler. As implementações
foram avaliadas quanto à complexidade algorítmica, legibilidade, capacidade de
abstração, precisão dos resultados e desempenho computacional, estabelecendo
métricas comparativas entre os paradigmas funcional e imperativo.

Os resultados evidenciam diferenças significativas na organização estrutural
do código, no tratamento de estados e efeitos colaterais, bem como na
expressividade das abstrações matemáticas. Observou-se que o paradigma
funcional favorece maior modularidade e aproximação conceitual com a formulação
matemática dos métodos, enquanto o paradigma imperativo apresenta vantagens
práticas relacionadas à familiaridade e controle explícito de fluxo. A
pesquisa contribui para a análise crítica da adequação de paradigmas de
programação no desenvolvimento de aplicações científicas e numéricas.

\noindent
\textbf{Palavras-chave:} Métodos Numéricos; Paradigmas de Programação;
Programação Funcional.

\subsubsection{Abstract}

\noindent
\textbf{Title:} Numerical Methods Implementation Comparison in Different
Programming Paradigms.

\vspace{0.5em}

This work investigates the implementation of numerical methods across
different programming paradigms, focusing on a comparison between a
multiparadigm language, Python, and a purely functional language, Haskell.
Classical numerical analysis methods were implemented, including the
Fixed-Point Method, the Newton--Raphson Method, and Euler's Method. The
implementations were evaluated in terms of algorithmic complexity, readability,
abstraction capability, result accuracy, and computational performance,
establishing comparative metrics between the functional and imperative
paradigms.

The results reveal significant differences in code structure, state management,
and handling of side effects, as well as in the expressiveness of mathematical
abstractions. The functional paradigm demonstrated greater modularity and
closer conceptual alignment with the mathematical formulation of the methods,
whereas the imperative paradigm showed practical advantages related to
familiarity and explicit flow control. This research contributes to a critical
analysis of the suitability of programming paradigms for scientific and
numerical computing applications.

\noindent
\textbf{Keywords:} Numerical Methods; Programming Paradigms; Functional
Programming.

\subsection{Introdução}

A crescente demanda por desempenho e confiabilidade em aplicações científicas
e computacionais tem ampliado o interesse na escolha adequada de paradigmas de
programação para implementação de algoritmos matemáticos. Métodos numéricos
constituem a base de diversas áreas da engenharia e da ciência da computação,
sendo amplamente utilizados na resolução de equações não lineares, modelagem
de sistemas dinâmicos e aproximação de soluções analíticas. Tradicionalmente,
tais métodos são implementados em linguagens de paradigma imperativo, devido
ao controle explícito de estado e fluxo de execução que esse modelo oferece.

Entretanto, o avanço das linguagens funcionais e multiparadigma tem levantado
questionamentos sobre sua adequação para aplicações numéricas. O paradigma
funcional, fundamentado em conceitos como funções puras, imutabilidade e
composição de funções, aproxima-se conceitualmente da formalização matemática
dos algoritmos, o que pode favorecer maior clareza e modularidade. Por outro
lado, há debates na literatura acerca de possíveis impactos no desempenho
computacional e na eficiência de memória quando comparado ao modelo imperativo
tradicional.

A literatura clássica de Análise Numérica, representada por autores como
Márcia Ruggiero, Steven C. Chapra, Richard L. Burden e J. Douglas Faires,
enfatiza aspectos como convergência, estabilidade, erro de truncamento e
eficiência computacional na implementação de métodos iterativos. Esses
critérios são centrais para avaliar não apenas o método matemático em si, mas
também sua realização computacional concreta.

Diante desse contexto, este trabalho tem como objetivo investigar a
implementação de métodos numéricos clássicos em dois paradigmas distintos,
utilizando Python, linguagem multiparadigma amplamente empregada em computação
científica, e Haskell, linguagem funcional pura. Busca-se analisar diferenças
estruturais, expressividade algorítmica, tratamento de estados e desempenho
computacional, contribuindo para uma compreensão mais aprofundada da adequação
de diferentes paradigmas no desenvolvimento de aplicações científicas e
numéricas.

\subsection{Metodologia}

O estudo foi conduzido em ambiente computacional controlado, utilizando um
computador com sistema operacional Windows 10, processador AMD Ryzen 5 3600X,
32 GB de memória RAM e GPU NVIDIA RTX 2060. As implementações em Python foram
desenvolvidas na versão 3.13, enquanto as implementações em Haskell foram
compiladas utilizando a versão mais recente do compilador GHC disponível em
2025.

No caso do Haskell, foram avaliados diferentes níveis de otimização do
compilador, variando desde ausência de otimização até o nível máximo
(\texttt{-O3}), com o objetivo de analisar o impacto das otimizações no
desempenho dos métodos numéricos. Em Python, foram utilizados recursos da
biblioteca padrão e bibliotecas científicas amplamente empregadas em computação
numérica, incluindo NumPy para operações vetoriais, Pandas para organização dos
dados experimentais e Matplotlib para visualização gráfica dos resultados. Em
Haskell, as implementações foram realizadas utilizando apenas a Prelude padrão,
sem dependência de bibliotecas matemáticas externas.

Foram implementados o Método do Ponto Fixo e o Método de Newton--Raphson, com
foco na análise de comportamento iterativo, convergência e desempenho
computacional. O Método de Euler foi desenvolvido, mas não constituiu o foco
central das análises comparativas. Para os métodos iterativos, foram adotados
critérios de parada baseados em tolerância numérica, utilizando
$10^{-6}$ como valor padrão e, em determinados experimentos, $10^{-3}$ para
avaliação do impacto na convergência. O número máximo de iterações foi definido,
na maioria dos testes, como 100 iterações, com variações controladas conforme a
necessidade experimental.

As implementações em ambas as linguagens mantiveram equivalência lógica,
utilizando as mesmas funções de teste, valores iniciais e critérios de parada,
garantindo comparabilidade direta entre os paradigmas funcional e imperativo.
O desempenho foi avaliado por meio da medição do tempo de execução sob as
mesmas condições de hardware e sistema operacional, considerando diferentes
níveis de otimização no caso do Haskell.

Foram realizadas múltiplas execuções para cada configuração experimental, com
posterior organização e análise dos dados em Python. Além do tempo de execução,
foram analisados aspectos estruturais do código, como legibilidade,
modularidade, organização de funções e tratamento de estados, por meio de
comparação qualitativa entre as implementações.

Os experimentos foram conduzidos sob condições equivalentes de carga do
sistema, buscando minimizar interferências externas. Os resultados foram
organizados em tabelas e gráficos para permitir análise comparativa consistente.

\subsection{Resultados e Discussões}

Os experimentos evidenciaram diferenças mensuráveis entre as implementações
realizadas em Python e Haskell, tanto no desempenho computacional quanto na
organização estrutural do código. Em relação ao tempo de execução, observou-se
que as versões compiladas em Haskell apresentaram desempenho significativamente
influenciado pelo nível de otimização do compilador. Sem otimizações, o
desempenho mostrou-se inferior ou comparável ao Python em determinados
cenários. Entretanto, com níveis elevados de otimização
(\texttt{-O2} e \texttt{-O3}), o tempo de execução apresentou redução
consistente, evidenciando o impacto direto da etapa de compilação na eficiência
final.

No caso do Python, o uso de bibliotecas como NumPy contribuiu para melhorar o
desempenho em operações numéricas vetorizadas, reduzindo parte da desvantagem
inerente à interpretação da linguagem. Ainda assim, a execução permaneceu
dependente do ambiente interpretado e do gerenciamento dinâmico de tipos, o que
impactou o tempo total de processamento em iterações mais intensivas.

Quanto à convergência dos métodos do Ponto Fixo e de Newton--Raphson, não foram
observadas diferenças matemáticas entre as linguagens, uma vez que os critérios
de parada, tolerâncias e funções testadas foram mantidos equivalentes. Os
resultados numéricos apresentaram precisão compatível com os valores esperados
na literatura clássica de Análise Numérica, conforme descrito por Steven C.
Chapra e Richard L. Burden, confirmando a correção algorítmica das
implementações.

Do ponto de vista estrutural, a implementação em Haskell demonstrou maior
modularidade e proximidade conceitual com a formulação matemática dos métodos.
A ausência de estado mutável e a utilização de funções puras favoreceram uma
organização mais declarativa. Esse resultado está alinhado com as discussões
sobre programação funcional apresentadas por John Hughes, que destaca a
clareza e composição funcional como vantagens estruturais do paradigma.

Por outro lado, a implementação em Python mostrou-se mais direta em termos de
fluxo de controle explícito e manipulação iterativa, especialmente para
usuários já familiarizados com o paradigma imperativo. A legibilidade, nesse
caso, esteve fortemente associada à experiência prévia do programador,
sugerindo que fatores humanos também influenciam a percepção de clareza do
código.

A análise comparativa indica que o paradigma funcional oferece vantagens
conceituais e estruturais, especialmente na representação formal de algoritmos
matemáticos. Contudo, o desempenho prático depende fortemente do modelo de
execução e das otimizações aplicadas. Já o paradigma imperativo, especialmente
quando combinado com bibliotecas otimizadas, mantém competitividade em
aplicações numéricas.

Esses resultados reforçam a ideia de que a escolha do paradigma não deve ser
ideológica, mas orientada pelo contexto da aplicação, requisitos de desempenho
e objetivos do projeto. A literatura de métodos numéricos enfatiza precisão e
estabilidade; este estudo acrescenta que a forma como o algoritmo é estruturado
também impacta a experiência de desenvolvimento e potencial manutenção futura.


\subsection{Conclusões}

Este estudo analisou a implementação de métodos numéricos clássicos sob dois
paradigmas de programação distintos, utilizando Python como linguagem
multiparadigma e Haskell como linguagem funcional pura. Os resultados
demonstraram que, do ponto de vista matemático, ambos os paradigmas são
igualmente capazes de produzir soluções corretas e precisas quando mantidos os
mesmos critérios de parada, tolerância e condições iniciais.

Entretanto, diferenças relevantes foram observadas na organização estrutural
do código e no desempenho computacional. O paradigma funcional evidenciou maior
proximidade conceitual com a formulação matemática dos algoritmos, favorecendo
modularidade e clareza declarativa. Por outro lado, o paradigma imperativo
mostrou-se competitivo em termos práticos, especialmente quando apoiado por
bibliotecas numéricas otimizadas. No caso do Haskell, verificou-se que o
desempenho depende significativamente do nível de otimização aplicado durante a
compilação.

Conclui-se que a escolha do paradigma de programação para aplicações numéricas
deve considerar não apenas desempenho bruto, mas também aspectos como
manutenção, legibilidade, abstração e contexto do projeto. Não há superioridade
absoluta entre paradigmas, mas adequação ao problema proposto.

Como continuidade, recomenda-se a ampliação do estudo para métodos numéricos de
maior complexidade, como técnicas de integração de ordem superior ou métodos
para sistemas de equações diferenciais, bem como a inclusão de métricas
adicionais, como consumo de memória e análise formal de complexidade. Também
seria relevante investigar o impacto de bibliotecas funcionais especializadas
em Haskell, comparando-as com frameworks científicos consolidados em Python.

\subsection{Referências}

RUGGIERO, Márcia A. G.; LOPES, Vera L. R.
\textit{Cálculo numérico: aspectos teóricos e computacionais}.
2. ed. São Paulo: Pearson Education do Brasil, 1996.

CHAPRA, Steven C.; CANALE, Raymond P.
\textit{Métodos numéricos para engenharia}.
7. ed. Porto Alegre: AMGH, 2016.

BURDEN, Richard L.; FAIRES, J. Douglas.
\textit{Análise numérica}.
9. ed. São Paulo: Cengage Learning, 2011.